% =====================================================================
% Fig 2 -- Avacchedaka triple as typed bucket; two version-typed
% claims that DO NOT collide
% =====================================================================
\begin{figure}[!tbp]
\centering
\resizebox{0.92\linewidth}{!}{%
\begin{tikzpicture}[
  font=\small,
  every node/.style={rounded corners=2pt},
  triple/.style ={draw=black!60,fill=blue!4,inner sep=4pt,minimum height=10mm,
                  text width=58mm,align=center},
  qual/.style   ={draw=black!60,fill=green!6,inner sep=2pt,minimum width=22mm,
                  minimum height=6mm,align=center,font=\footnotesize\ttfamily},
  cond/.style   ={draw=black!60,fill=orange!8,inner sep=2pt,minimum width=42mm,
                  minimum height=6mm,align=center,font=\footnotesize\ttfamily},
  prec/.style   ={draw=black!60,fill=violet!8,inner sep=2pt,minimum width=14mm,
                  minimum height=6mm,align=center,font=\footnotesize\ttfamily},
  bucket/.style ={draw=black!50,dashed,inner sep=4mm,rounded corners=4pt},
  arr/.style    ={-{Latex[length=2mm]},gray!70},
  >=Latex
]

% --- Generic schema box (top) ---
\node[triple] (schema) at (0,0)
  {\textbf{Avacchedaka triple}\\[1pt]
   $\langle\, a:\textit{qualificand},\;
     Q:\textit{qualifier},\;
     R:\textit{condition (limitor)} \,\rangle$
   \\[1pt] with precision $p\in[0,1]$};

% --- Two non-colliding claims (lower row) ---
\node[draw=black!50,fill=gray!2,inner sep=3mm,
      below left =14mm and -22mm of schema,
      label={[font=\footnotesize\bfseries]above:Claim A}] (A) {
  \begin{tikzpicture}[every node/.style={rounded corners=2pt}]
    \node[qual] (qA) at (0,0)   {Django.URLResolver};
    \node[qual,below=2mm of qA] (PA) {raises Resolver404};
    \node[cond,below=2mm of PA] (cA) {as of Django 4.2 LTS};
    \node[prec,right=2mm of cA] (pA) {p = 0.92};
  \end{tikzpicture}};

\node[draw=black!50,fill=gray!2,inner sep=3mm,
      below right=14mm and -22mm of schema,
      label={[font=\footnotesize\bfseries]above:Claim B}] (B) {
  \begin{tikzpicture}[every node/.style={rounded corners=2pt}]
    \node[qual] (qB) at (0,0)   {Django.URLResolver};
    \node[qual,below=2mm of qB] (PB) {raises NoReverseMatch};
    \node[cond,below=2mm of PB] (cB) {as of Django 5.0};
    \node[prec,right=2mm of cB] (pB) {p = 0.95};
  \end{tikzpicture}};

% --- Verdict ---
\node[draw=black!60,fill=green!10,rounded corners=3pt,
      inner sep=4pt,font=\small\bfseries,below=8mm of $(A)!0.5!(B)$] (ok)
  {No conflict: same qualificand, same qualifier-slot,
   \emph{different limitor (condition)}.};

% --- Arrows ---
\draw[arr] (schema) -- (A);
\draw[arr] (schema) -- (B);
\draw[arr,thick,green!50!black] (A) -- (ok);
\draw[arr,thick,green!50!black] (B) -- (ok);

\end{tikzpicture}}
\caption{\textbf{Avacchedaka as a typed bucket (Navya-Ny\=aya).}
Every \texttt{ContextItem} stamped into the store carries a four-tuple
$\langle a, Q, R, p\rangle$: the \emph{qualificand} $a$ (vi\'se\d{s}ya),
the \emph{qualifier} $Q$ (prak\=ara), the \emph{limitor / condition} $R$
(avacchedaka), and a precision $p\in[0,1]$.
Two claims that share $a$ but disagree on $Q$ \emph{do not} collide if
their $R$ differs --- here, ``raises \texttt{Resolver404}
\emph{as of Django 4.2 LTS}'' and ``raises \texttt{NoReverseMatch}
\emph{as of Django 5.0}'' coexist truthfully.
This is the operational definition the system uses to decide
\emph{when} a sublation is required (Section 4.3, Fig.~\ref{fig:sublation}).}
\label{fig:avacchedaka}
\end{figure}
